\documentclass[11pt,a4paper]{article}
\usepackage{fontspec,amsmath,amssymb,amsthm,microtype}
\setmainfont{lmroman10-regular.otf}[BoldFont=lmroman10-bold.otf,ItalicFont=lmroman10-italic.otf,BoldItalicFont=lmroman10-bolditalic.otf]
\usepackage{xcolor}
\usepackage[most]{tcolorbox}
\usepackage{mathrsfs}
\providecommand{\lt}{<}
\providecommand{\gt}{>}
\setmonofont{DejaVuSansMono.ttf}[BoldFont=DejaVuSansMono-Bold.ttf,ItalicFont=DejaVuSansMono-Oblique.ttf,BoldItalicFont=DejaVuSansMono-BoldOblique.ttf,Scale=0.86]
\newfontfamily\leanfallback{STIXTwoMath-Regular.otf}[Scale=0.86]
\definecolor{leanrule}{gray}{0.62}
\definecolor{leanback}{gray}{0.962}
\definecolor{leantitle}{gray}{0.22}
\definecolor{leanmuted}{gray}{0.38}
\newtcolorbox{leanbox}[1]{enhanced,breakable,colback=leanback,colframe=leanrule,
  boxrule=0.4pt,arc=0.6mm,left=2.4mm,right=2.4mm,top=1.8mm,bottom=1.8mm,
  before skip=0.8em,after skip=1em,fonttitle=\small\bfseries,coltitle=leantitle,
  colbacktitle=leanback,titlerule=0.3pt,title={#1}}
\newlength{\leancharwidth}
\newenvironment{leancode}{\par\smallskip\begingroup\ttfamily\footnotesize\raggedright
  \settowidth{\leancharwidth}{M}\setlength{\parskip}{0pt}\setlength{\parindent}{0pt}
  \frenchspacing}{\par\endgroup\smallskip}
\newcommand{\leanline}[2]{\hangindent=\dimexpr#1\leancharwidth+4\leancharwidth\relax
  \hangafter=1\noindent\hspace*{#1\leancharwidth}#2\par}
\newcommand{\leansource}[1]{\par{\footnotesize\color{leanmuted}#1}\par}
\newcommand{\leanlabel}[1]{\par\smallskip\noindent{\small\bfseries #1}\par}
\usepackage[margin=25mm]{geometry}
\usepackage[unicode,hidelinks]{hyperref}
\hypersetup{pdfauthor={LangevinLean Project}}
\setlength{\emergencystretch}{3em}
\setlength{\parskip}{0.4em}
\setlength{\parindent}{0pt}
\allowdisplaybreaks[2]
\raggedbottom
\title{\texorpdfstring{The first-order contraction and local-error inequality}{The first-order contraction and local-error inequality}}
\author{LangevinLean Project}
\date{}
\begin{document}\maketitle
Draft mathematical exposition; these notes may contain mistakes.


This theorem inserts an auxiliary exact kinetic path between the exponential numerical step and an exact reference state. Its proof combines the defined update's local error with continuous-time contraction. The reference velocity norm bound is an explicit hypothesis; establishing it from a stationary kinetic Gibbs law is proved separately.



The auxiliary exact path makes local numerical error depend on the current coupling discrepancy. Its velocity is controlled by a stationary Gaussian reference plus exact-flow contraction.


\section{\texorpdfstring{Auxiliary-trajectory velocity and the one-step discrepancy}{Auxiliary-trajectory velocity and the one-step discrepancy}}


Let \(A\) be the exact path starting at the numerical phase point and \(R\) the reference path, driven by the same additive input. Put \(e=\|A_0-R_0\|_T\). The integral-flow theorem gives \(\|A_s-R_s\|_T\le e^{-ms/\gamma}e\le e\), because \(m,s\ge0\). Velocity projection and the triangle inequality therefore give \(\|A_s^v\|\le\|A_s^v-R_s^v\|+\|R_s^v\|\le e+C\). The argument derives the exponential comparison from monotonicity of the exponential, multiplies it by the nonnegative initial norm, invokes the velocity projection bound, and combines the explicit reference bound with the triangle inequality.



The mild equation identifies the defined update's error relative to \(A_h\). Applying its proved local-norm theorem with \(C+e\) yields \(\|\operatorname{step}(A_0)-A_h\|_T\le3Mh^2(C+e)/4\). Finally insert \(A_h\) with the transformed triangle inequality and use exact-flow contraction between \(A_h\) and \(R_h\). The resulting inequality is

\[\|\operatorname{step}(A_0)-R_h\|_T\le e^{-mh/\gamma}e+\tfrac34Mh^2(C+e).\]

The update noise is exactly \((B(h),\int_0^hB(s)\,ds)\); the raw additive input and its velocity convolution are both specified. Thus the path equations and the numerical formula imply the one-step recurrence. Its application to random trajectories additionally requires the stated velocity moment bound and the specified Brownian input law.



Insert the auxiliary exact path between the numerical state and the stationary exact state. Exact-flow contraction controls its discrepancy from the reference. The triangle inequality bounds its velocity by the reference velocity radius plus that discrepancy, yielding a local error that depends linearly on the current coupling error.

\begin{leanbox}{Lean: \texttt{exponentialStep\_\allowbreak{}coupling\_\allowbreak{}bound}}
\leansource{\texttt{exponentialStep\_\allowbreak{}coupling\_\allowbreak{}bound} \quad KineticLangevin/\allowbreak{}OneStep.lean, lines 25–70}
\begin{leancode}
\leanline{0}{theorem\ exponentialStep\_coupling\_bound}
\leanline{4}{(f\ :\ E\ →\ ℝ)\ (m\ :\ ℝ)\ (M\ γ\ :\ ℝ≥0)}
\leanline{4}{(hdiff\ :\ Differentiable\ ℝ\ f)\ (hsc\ :\ StrongConvexOn\ univ\ m\ f)}
\leanline{4}{(hLip\ :\ LipschitzWith\ M\ (gradient\ f))\ (hmM\ :\ m\ ≤\ (M\ :\ ℝ))\ (hm\ :\ 0\ ≤\ m)}
\leanline{4}{(hγ\ :\ 0\ <\ (γ\ :\ ℝ))\ (hfriction\ :\ m\ +\ (M\ :\ ℝ)\ ≤\ (γ\ :\ ℝ)\ \textasciicircum{}\ 2)}
\leanline{4}{(A\ R\ W\ :\ ℝ\ →\ E\ ×\ E)\ (B\ :\ ℝ\ →\ E)}
\leanline{4}{(hA\ :\ Continuous\ A)\ (hR\ :\ Continuous\ R)\ (hB\ :\ Continuous\ B)}
\leanline{4}{(h\ C\ :\ ℝ)\ (hh\ :\ 0\ ≤\ h)\ (hC\ :\ 0\ ≤\ C)\ (hsmall\ :\ (γ\ :\ ℝ)\ *\ h\ ≤\ 1\ /\ 4)}
\leanline{4}{(heA\ :\ ∀\ t\ ∈\ Icc\ 0\ h,\ A\ t\ -\ A\ 0\ =}
\leanline{6}{-(∫\ s\ in\ 0..t,\ kineticNegativeDrift\ (gradient\ f)\ γ\ 1\ (A\ s))\ +\ W\ t)}
\leanline{4}{(heR\ :\ ∀\ t\ ∈\ Icc\ 0\ h,\ R\ t\ -\ R\ 0\ =}
\leanline{6}{-(∫\ s\ in\ 0..t,\ kineticNegativeDrift\ (gradient\ f)\ γ\ 1\ (R\ s))\ +\ W\ t)}
\leanline{4}{(hmild\ :\ ∀\ t\ ∈\ Icc\ 0\ h,\ (A\ t).1\ =\ psi\ γ\ 0\ t\ •\ (A\ 0).1\ -}
\leanline{6}{(∫\ s\ in\ 0..t,\ Real.exp\ (-(γ\ :\ ℝ)\ *\ (t\ -\ s))\ •\ gradient\ f\ (A\ s).2)\ +\ B\ t)}
\leanline{4}{(hkin\ :\ ∀\ t\ ∈\ Icc\ 0\ h,\ (A\ t).2\ -\ (A\ 0).2\ =\ ∫\ s\ in\ 0..t,\ (A\ s).1)}
\leanline{4}{(hreference\ :\ ∀\ t\ ∈\ Icc\ 0\ h,\ ‖(R\ t).1‖\ ≤\ C)\ :}
\leanline{4}{transformNorm\ γ\ 0}
\leanline{6}{(exponentialStep\ (gradient\ f)\ γ\ h\ (A\ 0)\ (B\ h,\ ∫\ s\ in\ 0..h,\ B\ s)\ -\ R\ h)\ ≤}
\leanline{6}{Real.exp\ ((-m\ /\ γ)\ *\ h)\ *\ transformNorm\ γ\ 0\ (A\ 0\ -\ R\ 0)\ +}
\leanline{8}{(3\ /\ 4\ :\ ℝ)\ *\ (M\ :\ ℝ)\ *\ h\ \textasciicircum{}\ 2\ *\ (C\ +\ transformNorm\ γ\ 0\ (A\ 0\ -\ R\ 0))}
\end{leancode}
\leanlabel{How the Lean proof follows the argument}
\texttt{exponential\allowbreak{}Step\_\allowbreak{}coupling\_\allowbreak{}bound} gives the formal counterpart of this step. The main intermediate facts are \texttt{e}, \texttt{he}, \texttt{hdec}, \texttt{hvelocity}, \texttt{hlocal}, \texttt{htriangle} in that order. Named proof applications include \texttt{Real.\allowbreak{}exp\_\allowbreak{}le\_\allowbreak{}one\_\allowbreak{}iff.\allowbreak{}mpr}, \texttt{neg\_\allowbreak{}nonpos.\allowbreak{}mpr}; the adjacent proof shows their hypotheses and order of use.
\end{leanbox}

\end{document}
